% Source: https://artofproblemsolving.com/community/c3h1606983_prove_abc_is_isosceles
% Stable identity: aops-topic-1606983-q01
Let $x=\angle PBC$ and $y=\angle PCB$. The angle sum of triangle $ABC$ is
\[
(10^\circ+40^\circ)+(20^\circ+x)+(30^\circ+y)=180^\circ,\qquad\text{so}\qquad x+y=80^\circ .
\]

\begin{asy}
usepackage("amsmath");
size(7cm);
defaultpen(fontsize(10pt));
real t=tan(radians(50));
pair A=(0,0),C=(6,0),B=(3,3*t);
pair P=(2.44562241,2.05212086);
draw(A--B--C--cycle,linewidth(0.9));
draw(A--P--C,blue);draw(B--P,blue);
dot(P);
label("$A$",A,SW);label("$B$",B,N);label("$C$",C,SE);label("$P$",P,dir(-75));
draw(arc(A,1.3,0,40),red);label("$40^\circ$",A+1.75*dir(20),red+fontsize(8pt));
draw(arc(A,1.9,40,50),red);pair a1=A+1.9*dir(45),a2=a1+(-0.95,0.45);draw(a1--a2,red+linewidth(0.3));label("$10^\circ$",a2,W,red+fontsize(8pt));
draw(arc(C,1.3,150,180),red);label("$30^\circ$",C+1.8*dir(165),red+fontsize(8pt));
draw(arc(B,1.3,degrees(A-B),degrees(P-B)),red);pair b1=B+1.3*dir((degrees(A-B)+degrees(P-B))/2),b2=b1+(-0.9,0.25);draw(b1--b2,red+linewidth(0.3));label("$20^\circ$",b2,W,red+fontsize(8pt));
draw(arc(B,0.7,degrees(P-B),degrees(C-B)),deepgreen);label("$x$",B+1.0*dir((degrees(P-B)+degrees(C-B))/2-8),deepgreen+fontsize(8pt));
draw(arc(C,0.75,degrees(P-C),degrees(B-C)),deepgreen);label("$y$",C+1.05*dir((degrees(P-C)+degrees(B-C))/2),deepgreen+fontsize(8pt));
label("\mbox{}",(min(currentpicture,user=true)+max(currentpicture,user=true))/2,black); // keeps the text after the figure black when LaTeX typesets the labels inline
\end{asy}

\textbf{Trigonometric Ceva.} For the concurrent cevians $AP$, $BP$, $CP$,
\[
\sin\angle PAB\cdot\sin\angle PBC\cdot\sin\angle PCA=\sin\angle PAC\cdot\sin\angle PBA\cdot\sin\angle PCB .
\]
(This follows from the law of sines in triangles $PAB$, $PBC$, $PCA$: for example $\dfrac{PA}{PB}=\dfrac{\sin\angle PBA}{\sin\angle PAB}$, and the product of the three ratios $\dfrac{PA}{PB}\cdot\dfrac{PB}{PC}\cdot\dfrac{PC}{PA}$ is $1$.) Here this reads
\[
\sin10^\circ\sin x\sin30^\circ=\sin40^\circ\sin20^\circ\sin(80^\circ-x),\qquad\text{i.e.}\qquad\dfrac{\sin x}{\sin(80^\circ-x)}=\dfrac{2\sin20^\circ\sin40^\circ}{\sin10^\circ}.
\]

\textbf{Evaluate the right side.} For any $\theta$, $\sin(60^\circ-\theta)\sin(60^\circ+\theta)=\sin^260^\circ-\sin^2\theta=\dfrac{3}{4}-\sin^2\theta$, so
\[
\sin\theta\sin(60^\circ-\theta)\sin(60^\circ+\theta)=\dfrac{3\sin\theta-4\sin^3\theta}{4}=\dfrac{\sin3\theta}{4}.
\]
With $\theta=20^\circ$: $\sin20^\circ\sin40^\circ\sin80^\circ=\dfrac{\sin60^\circ}{4}$. Using also $\sin80^\circ=\cos10^\circ$ and $2\sin10^\circ\cos10^\circ=\sin20^\circ$,
\[
\dfrac{2\sin20^\circ\sin40^\circ}{\sin10^\circ}=\dfrac{2\sin60^\circ}{4\sin80^\circ\sin10^\circ}=\dfrac{\sin60^\circ}{2\sin10^\circ\cos10^\circ}=\dfrac{\sin60^\circ}{\sin20^\circ}.
\]

\textbf{Solve for $x$.} So $\dfrac{\sin x}{\sin(80^\circ-x)}=\dfrac{\sin60^\circ}{\sin(80^\circ-60^\circ)}$. On $0^\circ<x<80^\circ$ the function $x\mapsto\dfrac{\sin x}{\sin(80^\circ-x)}$ is strictly increasing (the numerator increases and the positive denominator decreases), so it takes each value at most once. Hence $x=60^\circ$ and $y=20^\circ$.

\textbf{Conclusion.} $\angle A=50^\circ$, $\angle B=20^\circ+60^\circ=80^\circ$ and $\angle C=30^\circ+20^\circ=50^\circ$. Since $\angle A=\angle C$, the triangle is isosceles with $BA=BC$. $\blacksquare$
